\subsection{A six-element source and a one-row opposite corner}

\begin{lemma}[Proved: the positive six-element source shapes]
\label{CP2:A-six-positive-shapes}
Let $M$ have six elements and $H_2(BM;\mathbb Q)\ne0$.
Its minimal ideal is the $2\times2$ band
$K=\{e_{ij}:i,j\in\{0,1\}\}$, and
$M=K\amalg\{1,a\}$. Up to independent relabelling of the
rows and columns, the two actors of $a$ and its square are
exactly one of
\[
\begin{array}{c|c|c}
\sigma_a&\tau_a&a^2\\\hline
\mathrm{id}&\mathrm{id}&1\\
\mathrm{id}&\mathrm{id}&a\\
(01)&(01)&1\\
\mathrm{const}_0&\mathrm{const}_0&e_{00}\\
\mathrm{const}_0&\mathrm{const}_0&a.
\end{array}
\]
In each case $H_2(BM;\mathbb Q)$ is one-dimensional and
every rational degree-two square on $BM$ is zero.
\end{lemma}
\begin{proof}
The minimal-ideal and local-group reductions require at
least two rows and two columns. Six elements leave only
the displayed band and its two outside elements; a
nontrivial local group is too large. The two augmentation
modules are one-dimensional. Their balanced tensor is
nonzero exactly when the augmentation scalars of every
row and column actor agree. For a transformation of two
points the scalar is $1,-1,0$ for the identity, interchange
and a constant, respectively. Thus the actors of $a$ are
the three paired types in the table. The required actors
of $a^2$ then leave exactly the displayed products.
In the constant case the constant labels can independently
be normalized to zero. The balanced tensor has dimension
one, and Lemma~\ref{A:rank-one-square} gives the square
assertion.
\end{proof}

\begin{theorem}[Proved: a one-row opposite corner on at most four states]
\label{CP2:A-six-one-row-four-states}
Let $M$ be a six-element monoid, let $N$ be a finite monoid
whose minimal ideal has one row, and let $S$ be a nonempty
left-$N$, right-$M$ biset with at most four elements.
The associated two-object collage does not have
classifying space of type $\CP^2$.
\end{theorem}
\begin{proof}
Write $X$ for the collage space and $E=B(N,S,M)$ for its
two-sided action bar. The one-row hypothesis has two
consequences used separately below: $BN$ is rationally
acyclic, and the trivial \emph{left} $\mathbb QN$ module
is projective. For the latter, choose a minimal idempotent
$f$, with local group $G=fNf$. The average
$p=|G|^{-1}\sum_{g\in G}g$ obeys $np=p$ for every $n$,
because left multiplication permutes the group coordinate
in the unique row and the chosen column. Therefore
$\mathbb QN p\cong\mathbb Q$ is a projective summand.

If $H^2(BM;\mathbb Q)=0$, the relative-cup lemma already
excludes $X$. Otherwise use the preceding five source
types, and put $P=K^1\cong P_5$, $e=e_{00}$.
A state fixed by all band elements is fixed by $M$:
$sm=(se)m=s(em)=s$, since $em\in K$. Such a state makes
the source-corner inclusion nullhomotopic in $X$, so again
the relative-cup lemma applies.

Assume there are no fixed band states. The rectangular-band
act normal form then has regular two-point column cores.
If there are two cores, all four states are core states.
Write them $(b,j)$, where $j$ is the labelled column, and
put $h_{(b,j)}=e_{0j}$. For every $m\in M$ and $n\in N$,
\[
 h_s m=h_{sm},\qquad h_{ns}=h_s.
\]
These identities give the same equivariant natural
nullhomotopy of $E\to BM$ as in
Lemma~\ref{CP2:A-P5-full-regular-core}. Hence
$X\simeq BM\vee\operatorname{hocofib}(E\to BN)$.
The source summand already has one-dimensional second
cohomology and zero square. If $X$ were a projective plane,
its only second class would come from that summand and
would have zero square, a contradiction.

It remains to handle a unique core $\{c_0,c_1\}$.
Every band action on $S$ is then the constant map
$se_{ij}=c_j$, and every $N$ actor fixes $c_0,c_1$
individually. Write $A(s)=sa$.

\paragraph{Three split-algebra cases.}
For the first two source types, $a$ is central. In the
involution case put $p=(1+a)/2$; in the idempotent case
put $p=a$. The algebra $R=\mathbb QM$ splits as
\[
 R\cong\mathbb QP\times\mathbb Q,
 \qquad pR\cong\mathbb QP,
\]
and the trivial module lies in its $p$ summand. The monoid
map $\pi:M\to P$ fixes the band and sends $a$ to $1$.
Let $T$ be the actual quotient of $S$ by the $A$-orbits
in the involution case, or the image $AS$ in the idempotent
case. It is an $N$--$P$ biset, the quotient map is compatible
with $\pi$, and
$\mathbb Q[S]p\to\mathbb Q[T]$ is an isomorphism.

For the fifth source type, put $h=a-e$. Here $h$ is a
central idempotent, annihilates $\mathbb QK$ on both sides,
and $p=1-h$ gives the same algebra splitting, now with
$\pi(a)=e$. The set $AS$ is $N$ invariant, contains $c_0$
and does not contain $c_1$. Collapse $AS$ to $c_0$ to form
the actual $N$--$P$ quotient $T$. On permutation modules,
\[
 [s]p=[s]-[As]+[c_0].
\]
The vector $[c_0]$ and these vectors for $s\notin AS$
are a basis of $\mathbb Q[S]p$ and map to the evident
basis of $\mathbb Q[T]$.

In all three cases the resulting actual collage functor
induces a rational homology equivalence in every degree.
Indeed the trivial $M$ module is entirely in the $p$
block, and the displayed permutation-module isomorphism
identifies the derived right-action bar with that over
$P$. This comparison is $N$ linear, so it survives the
other derived tensor; the two corner maps and the
mapping-cylinder sequence give the assertion for the
collage. Its rational cohomology isomorphism respects
products because it is induced by an actual functor.
The target has at most four states. The proof of
Lemma~\ref{CP2:A-P5-four-state-acyclic-opposite} excludes
even a rational projective-plane cohomology ring: its
fixed-state, unique-core and full-core arguments use only
one-dimensional second cohomology and its nonzero square.
Thus these three cases are impossible.

\paragraph{The simultaneous-interchange case.}
Suppose $\sigma_a=\tau_a=(01)$ and $a^2=1$.
Set $v=e_{10}-e_{00}$ and $p_-=(1-a)/2$ in $R$.
An exact projective left resolution is
\[
 0\longrightarrow Re\xrightarrow{j}Rp_-
 \xrightarrow{d}Re\xrightarrow{\epsilon}\mathbb Q
 \longrightarrow0,
 \qquad d(p_-)=v,\qquad j(e_{i0})=e_{i0}p_-.
\]
The middle term has basis $p_-,e_{00}p_-,e_{10}p_-$;
$d$ has image $\mathbb Qv$ and kernel the last two
basis vectors, which are the displayed copy of $Re$.

The row interchange $\sigma(e_{ij})=e_{1-i,j}$,
$\sigma(a)=a$, fixes the right action on $S$ and gives
an automorphism $F$ of the whole category. Its semilinear
lift sends $e_{i0}$ in degree zero to $e_{1-i,0}$,
$rp_-$ in degree one to $-\sigma(r)p_-$, and $e_{i0}$
in degree two to $-e_{1-i,0}$. These maps commute with
both differentials. Tensoring with $\mathbb Q[S]$ gives
$\mathbb Qc_0$ in degrees zero and two, and
$\mathbb Q[S]p_-$ in degree one. The first differential
is zero, and the second sends $c_0$ to
$(c_0-c_1)/2$. This is an $N$-equivariant acyclic pair.
Its quotient is
\[
 \mathbb Q[0]\oplus T[1],
 \qquad T=\mathbb Q[S]p_-/\mathbb Q(c_0-c_1),
\]
with signs $+1,-1$ under $F$. Because $BN$ is rationally
acyclic, all positive rational homology of the bridge has
sign $-1$. The source resolution also gives homology only
in degrees zero and two, with sign $-1$ on $H_2(BM)$.
The mapping-cylinder sequence therefore gives sign $-1$
on $H_4(X)$.

If $X$ were a projective plane, a generator $\alpha$ in
degree two either restricts to zero on $BM$, when its
square vanishes by the relative-cup lemma, or has nonzero
restriction and hence $F^*\alpha=-\alpha$. Its square
then has sign $+1$, contrary to the sign on $H^4(X)$.

\paragraph{The nilpotent constant case: the exact relative complex.}
It remains to treat $\sigma_a=\tau_a=\mathrm{const}_0$
and $a^2=e$. Put
\[
 h=a-e,\qquad q=1-e,\qquad
 u_i=e_{i1}-e_{i0},\qquad v=e_{10}-e_{00}.
\]
Here $h$ is central, $h^2=0$, and $hK=Kh=0$.
The projectives $P_0=Re$ and $B=Rq$ have bases
$e_{00},e_{10}$ and $q,u_0,u_1,h$, respectively.
Let $j:Re\to B$ send $e_{i0}$ to $u_i$.
The following projective resolution continues in every
degree:
\[
 P_0=Re,\qquad P_1=B,\qquad P_k=Re\oplus B\quad(k\ge2),
\]
\[
 d_1(b)=bv,\qquad
 d_2(t,b)=j(t)+bh,\qquad
 d_k(t,b)=(0,j(t)+bh)\quad(k\ge3).
\]
To check exactness, $\ker\epsilon=\mathbb Qv$, while
$\ker d_1=\mathbb Q\{u_0,u_1,h\}$.
The map $j$ is an isomorphism onto the first two summands,
and right multiplication by $h$ sends $q$ to $h$ and
annihilates $u_0,u_1,h$. Thus the kernel of $d_2$, and
of each subsequent displayed map on its nonzero target
component, is exactly
$0\oplus\mathbb Q\{u_0,u_1,h\}$, the image of the
next differential. These assertions check exactness in
all degrees, not only through a finite prefix.

On the right permutation module put
\[
 D=\mathbb Q[S](1-e)
   =\operatorname{span}\{[s]-[c_0]:s\ne c_0\},
 \qquad d=[c_1]-[c_0],
 \qquad H(z)=zh.
\]
Then $H([s]-[c_0])=[As]-[c_0]$, $H^2=0$, and $Hd=0$.
Tensoring the resolution gives the $N$ complex
\[
 C_0=\mathbb Q,\quad C_1=D,\quad C_k=\mathbb Q\oplus D
 \ (k\ge2),
\]
with $d_1=0$, $d_2(t,z)=td+Hz$, and
$d_k(t,z)=(0,td+Hz)$ thereafter.

The augmentation $\mathbb Q[S]\to\mathbb Q$ models the
actual bridge projection to $BM$. On this resolution its
target complex is $\mathbb Q$ in degree zero and every
degree at least two, zero in degree one, and all target
differentials are zero. The chain map is the identity in
degree zero and the first-coordinate projection in all
degrees at least two. Its exact kernel is the $N$ complex
\[
 K_0=0,\qquad K_k=D\quad(k\ge1),\qquad
 d_K=H\quad(k\ge2).                                      \tag{*}
\]
Every target term is a trivial $N$ module. Since $BN$ is
rationally acyclic, its derived $N$ coinvariants compare
quasi-isomorphically to the target itself. Also collapsing
$BN$ in the collage is a rational homology equivalence.
The resulting cofiber is the cofiber of the actual
bridge projection to $BM$. The short exact kernel sequence
therefore gives
\[
 \widetilde C_*(X;\mathbb Q)\simeq
 \bigl(\mathbb Q\otimes_{\mathbb QN}^{\mathbf L}K\bigr)[1],
 \qquad
 H_2(X;\mathbb Q)\cong(D/HD)_N.                         \tag{**}
\]
The second assertion is the degree-one homology of the
derived tensor of a complex beginning in degree one.
This calculation uses the actual relative complex; it
does not presume that a vertical-homology summand survives
a spectral sequence.

The set map $A$ sends both core points to $c_0$ and has
$A^2$ equal to the constant $c_0$. With at most two other
states, its possibilities divide into the following three
exhaustive cases.

If $H=0$, the complex $(*)$ is
$\bigoplus_{k\ge1}D[k]$ with zero differential.
Every nonzero $\operatorname{Tor}^{\mathbb QN}_j(\mathbb Q,D)$
would consequently occur in arbitrarily high homological
degrees of $(**)$. Bounded homology forces all these Tor
groups to vanish, including coinvariants, and then
$H_2(X;\mathbb Q)=0$.

If $HD=\mathbb Qd=L$, this is a trivial left $N$ module
and hence projective. The surjection $H:D\to L$ has an
$N$-linear section with image $J\cong L$, so
$D=\ker H\oplus J$. Put $V=\ker H/L$.
There is an acyclic subcomplex of $(*)$ whose degree-one
term is $L$, whose terms in all higher degrees are
$L\oplus J$, and whose differential is the isomorphism
$J\to L$ on consecutive terms. The quotient is the
zero-differential complex
\[
 (D/L)[1]\ \oplus\ \bigoplus_{k\ge2}V[k],
 \qquad D/L\cong V\oplus J.
\]
Boundedness of $(**)$ forces every Tor group with $V$
to vanish because its shifts repeat indefinitely.
The remaining $J\cong\mathbb Q$ is projective on the
left, so it contributes only in degree two of $X$.
In particular $H_4(X;\mathbb Q)=0$.

In the last case $S=\{c_0,c_1,w,x\}$, after relabelling,
and
\[
 Ax=w,\qquad Aw=Ac_0=Ac_1=c_0,
 \qquad HD=\mathbb Q([w]-[c_0]).
\]
Every transformation commuting with $A$ and fixing the
two core points is determined by its value on $x$.
Its four possibilities are
\[
\begin{array}{c|cc}
 &x&w\\\hline
 I&x&w\\
 F&w&c_0\\
 P_0&c_0&c_0\\
 P_1&c_1&c_0.
\end{array}
\]
All fix $c_0,c_1$. With ordinary function composition,
$F^2=P_0$, $FP_i=P_i$, $P_iF=P_0$, and $P_iP_j=P_j$.
If both $P_0$ and $P_1$ occur in the image of $N$, the
module $D/HD$ has basis $d,[x]-[c_0]$ and its coinvariant
relations kill first $[x]-[c_0]$ using $P_0$, then $d$
using $P_1$. Formula $(**)$ gives $H_2(X;\mathbb Q)=0$.

If neither of these two idempotents occurs, only $I$
occurs, since $F^2=P_0$. The $N$ action is trivial.
Collapsing the rationally acyclic $BN$ gives an actual
functor to the ordinary $M$ act cone which is a rational
homology and cohomology-ring equivalence. Its homology
would be bounded, contrary to the bounded-small-act
theorem, for example
Theorem~\ref{CP2:A-bounded-eleven-acts}.

If exactly one of $P_0,P_1$ occurs, it is central in the
whole image of $N$. Indeed the only extra possible actor
is $F$, and it can occur only with $P_0$, with which it
commutes. The corresponding transformation is therefore
a proper $N$--$M$ biset idempotent with image
$\{c_0,c_1\}$. The actual cut-retraction theorem
\ref{CP2:A-cut-retract} would preserve projective-plane
type in that two-cross-arrow subcollage, contradicting
Lemma~\ref{CP2:A-high-two-cross}.
This exhausts the last case and proves the theorem.
\end{proof}

\begin{corollary}[Proved: the opposite-order-eight scope for $(6,8,4)$]
\label{CP2:A-six-eight-four-target-scope}
For the eighteen-arrow profile $(6,8,4)$ with a positive
source corner, it suffices to test the five sources of
Lemma~\ref{CP2:A-six-positive-shapes} against all $1205$
ordinary order-eight monoids having at least two minimal
rows and two minimal columns. Every table in that target
catalog is required, including those with zero rational
second homology.
\end{corollary}
\begin{proof}
A one-row opposite ideal is excluded by the theorem.
A one-column ideal gives exact projective-right
coinvariants, hence an actual orbit quotient inducing a
rational cohomology-ring isomorphism onto an ordinary
six-element-monoid act cone. Its bounded rational homology
forces zero products by the bounded-small-act theorem.
The remaining order-eight ideals have trivial group and
shapes $2\times2$, $2\times3$, or $3\times2$.
Their complete ordinary catalogs have
$1181+12+12=1205$ tables, as already proved and checked in
Proposition~\ref{CP2:A-P5-eight-corner-scope}.
This corollary closes the input reduction only; it makes
no claim that the new four-state finite biset search has
been completed.
\end{proof}
